% !TeX root = main.tex
\chapter*{АНОТАЦІЯ}
\addcontentsline{toc}{chapter}{АНОТАЦІЯ}

Денисюк Д.~О. Інформаційна технологія виявлення аномальних кодів у фото- та відеооб’єктах веб-ресурсів. Кваліфікаційна наукова праця на правах рукопису.

Дисертація на здобуття наукового ступеня доктора філософії за спеціальністю 126 «Інформаційні системи та технології» галузі знань 12 «Інформаційні технології». Хмельницький національний університет, Хмельницький, 2026.

Практична проблема дослідження полягає в тому, що фото- та відеооб’єкти надходять до веб-ресурсів через штатні канали, але на подальшому маршруті їх перевіряють, декодують, перетворюють і передають компоненти з різними правилами визначення типу, меж контейнера та допустимих операцій. Тому формальна коректність і візуальна правдоподібність вхідного об’єкта не виключають слабких просторових або частотних змін, нетипових службових чи додаткових байтових областей і подій, що виникають лише під час оброблення. Окремі структурні, спектрально-зорові та подієво-поведінкові засоби спостерігають лише частину цього маршруту. Наукова проблема полягає у формуванні обґрунтованого рішення за неповних, різнорідних і залежних від маршруту спостережень. Суперечність виявляється в тому, що вибір найбільшої локальної оцінки робить рішення чутливим до випадкового сильного відхилення, тоді як вимога одночасного підтвердження всіма каналами спричиняє пропуск за недоступної модальності або неактивованого механізму.

Доказовий зміст результату обмежено розмежуванням основних понять. Аномалія є спостережуваним відхиленням і сама по собі не доводить загрози. Аномальним кодом названо цілеспрямовано сформовану послідовність даних, команду, фрагмент коду або структурне значення, що не відповідає заявленому призначенню носія та може підтримати несанкціоновану дію. Прихована загроза є сценарієм можливого небезпечного впливу, для якого встановлюється зв’язок між властивостями носія, операцією оброблення та очікуваним наслідком. Про шкідливий код ідеться лише за підтверджених програмної семантики та шкідливої функції.

Дослідницьке припущення полягає в тому, що узгодження ознак за спільним ідентифікатором об’єкта, походженням, локалізацією, причинним контекстом і маскою доступності дає змогу формувати простежуване рішення без прирівнювання відсутньої модальності до нормального стану. Метою роботи є формування інформаційної технології виявлення аномальних кодів у фото- та відеооб’єктах веб-ресурсів на основі однієї моделі та двох взаємопов’язаних методів. Об’єктом дослідження є процес такого виявлення під час приймання й оброблення мультимедійних об’єктів компонентами веб-ресурсу. Предмет дослідження становлять модель мультимодального представлення прихованої загрози, метод спектрально-зорового та структурного аналізу, метод поведінкового узгодження ознак різного походження та інформаційна технологія їх послідовного застосування.

Розроблено модель мультимодального представлення прихованої загрози. У ній мультимедійний об’єкт подано як багаторівневий інформаційний контейнер, для якого незмінне байтове подання, контрольовано декодований вміст, структурна карта, технічний контекст і причинно пов’язаний журнал подій залишаються різними, але простежувано пов’язаними сутностями. Структурні, спектрально-зорові та подієво-поведінкові ознаки інтегруються у представлення $Z$ зі збереженням часткових оцінок, локалізацій, походження та доступності. Оцінку ризику, упевненість, аналітичний стан, пояснення й технологічну дію розмежовано, тому неповне спостереження не перетворюється на автоматичний висновок про норму.

Удосконалено метод спектрально-зорового та структурного аналізу завдяки узгодженню двох статичних гілок зі збереженням їхніх окремих оцінок і локалізацій. Спектрально-зорова гілка формує високочастотні залишки, характеристики дискретного косинусного й дискретного хвилькового перетворень, ентропійні ознаки та часову агрегацію кадрів. Структурна гілка перевіряє заголовки, межі елементів контейнера, службові поля, завершальні маркери й додаткові області. Емпірична перевірка детермінованих складників охопила PNG, JPEG, WebP, MP4 і WebM. Параметри первинно визначено для PNG і MP4, а перенесення на JPEG, WebP і WebM оцінено без повторного навчання або добору порогів. Повну модель зорового перетворювача не навчено й окремо не перевірено.

Запропоновано метод поведінкового узгодження ознак різного походження, за яким статичний опис зіставляється лише з подіями оброблення того самого об’єкта, що мають підтверджені ідентифікатор, часові межі, фазу та причинний зв’язок. Метод передбачає кодування подій, фазову маску, керовану інтеграцію, оцінювання повноти й посилання на незмінний журнал для контрольної перевірки. Відсутність придатного журналу позначається як недоступність модальності. Для експериментального оцінювання реалізовано скорочений форматно-адаптивний навчуваний шлюз, який узгоджує структурну, спектрально-зорову та спільну статичні оцінки, їхні розбіжності й відомості про формат. Він не опрацьовує причинно пов’язаних послідовностей подій, а технічне поведінкове спрацювання використовує лише як окрему запобіжну умову. Тому результати шлюзу характеризують скорочене узгодження статичних оцінок і не є перевіркою повного подієво-поведінкового кодувальника та перехресного зіставлення.

Набула подальшого розвитку інформаційна технологія, у якій модель і два методи організовано в п’ять етапів: попереднє оброблення, виділення ознак, інтеграцію ознак, прийняття рішення та реєстрацію результатів. Початковий об’єкт зберігається незмінним, реконструйований кандидат проходить окремий цикл, а типізований вихідний стан містить версії складників, локалізації, маску доступності та посилання на журнал. Аналітичний результат відокремлено від дії політики: допуску, експертної перевірки, реконструкції або блокування. Практичне значення полягає у визначенні простежуваного контуру між штатним каналом приймання мультимедійного об’єкта та прикладною логікою веб-системи, а також складу машинозчитуваного запису для повторної перевірки. Для попереднього контрольного опису повного пакета немає й повторного локального запуску тієї конфігурації не виконано. Для внутрішньої та чотирьох зовнішніх серій збережено пооб’єктні реєстри, програмний код, конфігурацію, ознаки, рішення, часові записи, параметри моделей, відомості про середовище та контрольні суми.

Перший цикл фактичної перевірки подано як описово-статистичне та арифметичне оцінювання контрольного опису. Він характеризує набір зі $100$ PNG і $10$ коротких MP4, до якого входили $55$ контрольних і $55$ навмисно змінених об’єктів. Позитивна еталонна мітка позначала контрольовано внесене структурне або статистичне відхилення, а не доведену загрозу чи виконання шкідливого коду; контрольні рядки не були виконуваним програмним кодом. Похідні різних категорій могли походити від тих самих первинних носіїв, групового поділу до створення похідних не виконано, а параметри налаштовували й оцінювали на тому самому наборі.

З агрегованої матриці рішень із $50$ істинно позитивними ($TP$), $48$ істинно негативними ($TN$), $7$ хибнопозитивними ($FP$) і $5$ хибнонегативними ($FN$) рішеннями арифметично одержано частку правильних рішень $0{,}8909$, точність позитивного виявлення $0{,}8772$, повноту $0{,}9091$, $F_1=0{,}8929$, частку хибнопозитивних рішень $0{,}1273$ і частку хибнонегативних рішень $0{,}0909$. Значення площі під кривою залежності повноти від частки хибнопозитивних рішень $0{,}9367$ лише перенесено з контрольного опису: без неперервних пооб’єктних оцінок його незалежне відтворення неможливе. У внутрішньому компонентному зіставленні на тому самому наборі вилучення структурних ознак зменшило повноту на $40{,}00$ відсоткового пункту, а вилучення спектрально-зорових ознак --- на $34{,}55$ відсоткового пункту. Вилучення скороченого правилового подієво-поведінкового вектора не змінило агрегованих показників. Це зіставлення виконано без повторного навчання, парної статистичної перевірки та коректного зовнішнього порівняння.

Сума середніх значень чотирьох зафіксованих у звіті операцій становить $418{,}09$~мс, що арифметично відповідає $2{,}39$ об’єкта за секунду в послідовному режимі. На спектрально-зорову операцію припадає $98{,}49$~\% цього часу. Показник не охоплює окремо прийняття рішення й реєстрацію результату, тому не характеризує наскрізний час п’яти етапів або пропускну здатність розгорнутої веб-системи. Зважене середнє категоріальних часових даних дорівнює $424{,}72$~мс і не узгоджується з операційним зведенням часу; ці категоріальні значення не використано для підсумкового оцінювання.

У внутрішній групово відокремленій серії детерміновано сформовано $200$ первинних груп PNG і MP4, кожна з яких містила контрольний і контрольовано змінений об’єкти, тобто $400$ файлів. Групи розподілено на навчальну, налаштувальну й перевірну частини у співвідношенні $60:20:20$ до створення похідних. На недоторканій перевірній частині з $80$ об’єктів спільна конфігурація дала $40$ правильно визначених контрольних, жодного хибнопозитивного, два хибнонегативні та $38$ правильно визначених змінених об’єктів. Частка правильних рішень становила $0{,}9750$, повнота --- $0{,}9500$, $F_1=0{,}9744$, а площа під характеристичною кривою --- $0{,}9869$. Середній час формування ознак становив $53{,}89$~мс для PNG і $198{,}76$~мс для MP4.

Експериментальна перевірка охопила шість основних серій: ретроспективну, внутрішню групово відокремлену та чотири зовнішні. У зовнішніх серіях використано $38$ природних вихідних відеозаписів із $12$ пристроїв, об’єднаних у $24$ групи походження ``пристрій--сцена''. Перша зовнішня серія на $56$ похідних MP4 не відтворила внутрішньої переваги спільного опису над структурною конфігурацією: їх збалансовані частки правильних рішень становили відповідно $0{,}7857$ і $0{,}8333$, тоді як для спектрально-зорової конфігурації одержано $0{,}5000$. Наступна серія на $32$ об’єктах показала, що вимога додаткового підтвердження усуває чотири хибнопозитивні рішення, але збільшує кількість пропусків від чотирьох до восьми.

Великомасштабна серія містила $10\,000$ фізичних об’єктів: $500$ базових відеофрагментів було подано у форматах PNG, JPEG, WebP, MP4 і WebM за чотирьох сценаріїв. До цієї кількості входять раніше зафіксовані $2\,000$ MP4 та $8\,000$ нових форматних похідних. Тому $10\,000$ файлів не є $10\,000$ незалежними джерелами: вони походять із восьми записів і чотирьох груп ``пристрій--сцена''. На цьому матеріалі збалансована частка правильних рішень становила $0{,}7636$ для структурної, $0{,}5001$ для спектрально-зорової, $0{,}7916$ для спільної конфігурації та $0{,}8333$ для скороченого форматно-адаптивного навчуваного шлюзу. Останнє значення належить циклу розроблення й потребувало окремої перевірки на нових пристроях.

Незалежну перевірку шлюзу виконано на $2\,000$ фізичних об’єктах, утворених зі $100$ базових фрагментів восьми нових записів чотирьох пристроїв. Після фіксації моделі й порога збалансована частка правильних рішень змінилася від $0{,}7827$ для спільної конфігурації до $0{,}8307$ для шлюзу. Він усунув $76$ хибнопозитивних рішень, але додав $84$ хибнонегативні й пропустив усі $500$ об’єктів лише зі спектральною зміною. Точний двобічний критерій Мак-Немара для звичайної частки правильних рішень дав $p=0{,}5801$. Отже, встановлено відтворюваний компроміс між специфічністю та повнотою, а не безумовну перевагу шлюзу за всіма показниками.

Окреме ресурсне випробування виконано на реальних файлах у п’яти незалежних процесних запусках. Для структурної, спільної конфігурацій і навчуваного шлюзу 95-й процентиль часу оброблення становив відповідно $2{,}136$, $853{,}971$ і $1078{,}904$~мс, а медіана найбільшого обсягу фізичної пам’яті процесу --- $78{,}77$, $101{,}53$ і $110{,}05$~МіБ. Розвідувальний інтегральний показник $K_E$, що поєднує достовірність виявлення, часові витрати, пам’ять і стійкість за форматами, дорівнював відповідно $77{,}50$, $68{,}95$ і $69{,}46$ бала. Структурна конфігурація мала найбільше значення $K_E$, але не пройшла форматного допуску через результат для WebM; спільна конфігурація та шлюз цей допуск пройшли. Показник $K_E$ є вторинним рейтингом для наявного середовища, а не відсотком правильних рішень або підтвердженням промислової готовності.

Межі висновків визначаються залежністю форматних і сценарних похідних від спільних базових фрагментів та невеликою кількістю незалежних пристроїв. Причинно пов’язаних журналів подій і повного представлення $Z_b$ не сформовано, а у $4\,832$ обробленнях із доступною технічною поведінковою ознакою не зафіксовано жодного спрацювання. Це не є перевіркою повного методу поведінкового узгодження. Не оцінено виконуваний шкідливий код, паралельне навантаження та промисловий п’ятиетапний контур, тому результати не обґрунтовують безумовного автоматичного блокування.

За темою дисертації опубліковано дев’ять наукових праць: дві статті містять основні наукові результати й опубліковані у фахових наукових виданнях України, сім праць засвідчують апробацію матеріалів у виданнях IEEE, CEUR Workshop Proceedings і матеріалах наукової конференції. Повні бібліографічні описи наведено у списку публікацій здобувача.

\textbf{Ключові слова:} аномальний код, мультимедійний об’єкт, структурний аналіз, спектрально-зоровий аналіз, інформаційна технологія, обчислювальна ефективність.

\clearpage
